\documentclass[11pt]{article}
\usepackage{fontspec,amsmath,amssymb,amsthm,hyperref}
\usepackage[margin=1in]{geometry}
\setlength{\emergencystretch}{3em}
\title{Quadratic sharpness and the singular critical boundary}
\author{LangevinLean mathematical exposition}
\date{}
\begin{document}\maketitle
\noindent Source module: \texttt{KineticLangevin/QuadraticSharpness.lean}. Compile this proof document with XeLaTeX or LuaLaTeX. The accompanying HTML pairs the mathematical argument with the exact Lean statements and records the scope of the formalization.


Exact solutions of the quadratic kinetic equations justify the slow exponential rate and explain why the repeated-root boundary cannot have the singular prefactor substituted into it. Both examples start with unit position difference and equal initial velocities.


\section{Construct the genuine overdamped quadratic response}


Let \(0\le a<b\), set friction \(\gamma=a+b\), and curvature \(m=ab\). The actual deterministic difference solution is \[x(t)=\frac{be^{-at}-ae^{-bt}}{b-a},\qquad v(t)=\frac{ab(e^{-bt}-e^{-at})}{b-a}.\] Lean checks \(x(0)=1\), \(v(0)=0\), then differentiates both exponentials and proves exactly \(x\prime=v\), \(v\prime=-\gamma v-mx\). These are the synchronous kinetic difference equations for the quadratic potential \(mx^2/2\); shared Brownian forcing cancels. Thus the example respects the equal-initial-velocity premise of the source contraction statement.



\noindent Lean declarations: \texttt{quadraticPositionResponse}, \texttt{quadraticVelocityResponse}, \texttt{quadraticPositionResponse\_zero}, \texttt{quadraticVelocityResponse\_zero}, \texttt{quadraticResponse\_derivatives}.

\section{Prove that no faster exponential rate is possible}


The position response satisfies \(x(t)\ge e^{-at}\) for all \(t\ge0\), since \(x(t)-e^{-at}=a(e^{-at}-e^{-bt})/(b-a)\ge0\). For any proposed faster exponent \(\beta>a\) and any finite real constant \(C\), choose \(t=\log(|C|+1)/(\beta-a)\). Then \(e^{-at}=(|C|+1)e^{-\beta t}>Ce^{-\beta t}\). Lean constructs this nonnegative time explicitly and proves the strict violation, so the sharpness assertion involves no informal asymptotics.



\noindent Lean declarations: \texttt{quadraticPositionResponse\_slow\_lower}, \texttt{quadraticPositionResponse\_no\_faster\_rate}.

\section{Check the repeated-root solution and its unavoidable prefactor}


At critical damping, \(\gamma=2a\) and \(m=a^2\), the exact equal-initial-velocity solution is \(x(t)=(1+at)e^{-at}\), \(v(t)=-a^2te^{-at}\). Lean verifies the same kinetic equations directly. When \(a>0\), no finite constant can bound this position by \(Ce^{-at}\): choose \(t=(|C|+1)/a\). Thus the critical exponent is a limiting rate whose polynomial prefactor cannot be suppressed. This explains the singular denominator at the equal-curvature boundary of the displayed contraction prefactor.



\noindent Lean declarations: \texttt{criticalQuadraticPosition}, \texttt{criticalQuadraticVelocity}, \texttt{criticalQuadratic\_derivatives}, \texttt{criticalQuadratic\_prefactor\_obstruction}.

\section{Identify the source overdamped rate with the genuine quadratic roots}


For \(m>0\), \(\gamma>0\), and \(\gamma^2>4m\), put \(a=(\gamma-\sqrt{\gamma^2-4m})/2\), \(b=(\gamma+\sqrt{\gamma^2-4m})/2\). Lean proves \(0<a<b\), \(a+b=\gamma\), and \(ab=m\), using the exact square-root identity. Therefore the lower-bound example has precisely the slow rate reported in the source overdamped regime; at \(\gamma=2\sqrt M\) it gives \(a=\sqrt M-\sqrt{M-m}\).



\noindent Lean declarations: \texttt{quadratic\_slow\_fast\_roots}.

\end{document}
